%/*% pdflatex -file-line-error -interaction nonstopmode %

\documentclass{article}
\usepackage{graphicx}
\usepackage{amsmath}
\usepackage{amsfonts}
\usepackage{bookmark}
\usepackage{exercises}
\usepackage{amsthm}
\usepackage{yhmath}
\newcommand*{\oldneg}{\mathord{\sim}}

\theoremstyle{definition}
\newtheorem{definition}{Definition}

\setlength{\parindent}{.5in}
\setlength{\parskip}{12pt}

\title{1.1 Exercises}
\begin{document}
\maketitle

\section*{1}
Parabola.

\section*{2}
Given a quadratic function in the form $y=ax^2+bx+c$, if $a>0$ then the $y$-coordinate of the vertex is the minimum value, and if $a<0$, it is the maximum value.

\section*{3}
$y=4x^2$\newline
\begin{table}[h!]
	\begin{center}
		\begin{tabular}{l|c|c|c|c|r}
			$x$ & $-2$ & $-1$ & $0$ & $1$ & $2$ \\
			\hline
			$y$ & $16$ & $4$ & $0$ & $4$ & $16$
		\end{tabular}
	\end{center}
\end{table}
	    

\section*{19}
The $x$-coordinate of the vertex is given by $-\frac{b}{2a}$, so the $x$-coordinate acutally should be $-\frac{24}{2(4)}=-3$.

\section*{20}
The value of $c$ is $-7$, not $7$.

\section*{33}
Because $a$ is negative, there is a maximum value. The $x$-coordinate of the vertex is $-\frac{0}{2(-6)}=0$. So, the $y$-coordinate of the vertex is $-6(0)^2-1=-1$, which is 
also the maximum value of this function.

\section*{35}
Because $a$ is positive, there is a minimum value. The $x$-coordinate of the vertex is $-\frac{8}{2(2)}=-2$. So, the $y$-coordinate of the vertex is $2(-2)^2+8(-2)+7=-1$, 
which is also the minimum value of this function.

\section*{37}
Because $a$ is positive, there is a minimum value. The $x$-coordinate of the vertex is $-\frac{6}{2(\frac{3}{2})}=-2$. Then, the $y$-coordinate of the vertex is 
$\frac{3}{2}(-2)^2+6(-2)+4=-2$, which is also the minimum value of this function.

\section*{38}
Because $a$ is negative, there is a maximum value. The $x$-coordinate of the vertex is $-\frac{-7}{-2(\frac{1}{4})}=-14$. Then, the $y$-coordinate of the vertex is 
$-\frac{1}{4}(-14)^2-7(-14)+2=51$, which is also the maximum value of this function.

\section*{39}
D

\section*{40}
C

\section*{41}
$a=-0.02,\ b=1,\ c=6$

\section*{42}
$a=-0.01,\ b=0.7,\ c=6$

\section*{43}
$y=x^2-8x+3$, $y=x^2-8x+10$, $y=x^2-8x-12$

\section*{44}
C

\section*{45}
A

\section*{46}
B

\section*{53}
In a parabola, points on the parabola with the same $y$-coordinate are equidistant from the axis of symmetry. We are given two points $(2,3),\ (-4,3)$, which both have the 
same $y$-coordinates. So, they lie equidistant from the axis of symmetry. The midpoint of the line formed between these two points is $(-1,\ 3)$, and hence, the equation of 
the axis of symmetry is $x=-1$.

\section*{54}
We are given a function $y=ax^2+bx+c$. The $x$-coordinate of the vertex of this function is $-\frac{b}{2a}$. Then, the $y$-coordinate of the vertex is:\newline
$a(-\frac{b}{2a})^2+b(-\frac{b}{2a})+c$\newline$\rightarrow a(\frac{b^2}{4a^2})-\frac{b^2}{2a}+c$\newline$\rightarrow \frac{b^2}{4a}-\frac{b^2}{2a}+c$\newline$\rightarrow
\frac{b^2(2a)}{4a(2a)}-\frac{b^2(4a)}{2a(4a)}+c$\newline$\rightarrow \frac{b^2(2a)-b^2(4a)}{8a}+c$\newline$\rightarrow \frac{b^2(-2a)}{8a}+c$\newline$\rightarrow -\frac{b^2}
{4a}+c$

\section*{55}
This can be solved by finding the vertex of the function. First, we simplify the function:\newline $y=4000-80x+200x-4x^2\newline\rightarrow y=-4x^2+120x+4000$\newline
The $x$-coordinate of the vertex is $-\frac{120}{2(-4)}=15$. Then, the $y$-coordinate of the vertex is $-4(15)^2+120(15)+4000=4900$. So, the maximum revenue comes if the 
price is increased by $\$0.75$, leading to a total sale of $\$4,900$.

\section*{56}
We can represent this situation with the model for $R(x)$ the monthly revenue:\newline
$R(x)=(320-20x)\times (70+5x)\newline\rightarrow=22400+1600x-1400x-100x^2\newline\rightarrow =-100x^2-200x+22400$ Then, the $x$-coordinate of the vertex is $-\frac{-200}
{2(-100)}=-1$. Thus, the $y$-coordinate of the vertex is $22500$. Hence, the shop should reduce the price by $\$20$ to produce maximum monthly revenue.

\section*{57}
This can be solved by finding the minimum value, that is, the $y$-coordinate of the vertex. The $x$-coordinate of the vertex is:\newline
$-\frac{-\frac{7}{15}}{2(\frac{1}{9000}})=2100$, and hence, the $y$-coordinate of the vertex is $\frac{1}{9000}(2100)^2-\frac{7}{15}(2100)+500=10$. Thus, at it lowest point, 
the cable is $10$ ft. above the road.

\section*{58}
If it could, the maximum value of this jump would be greater than $3$. Well, the $x$-coordinate of the vertex of this function is $-\frac{1.3}{2(-0.2)}=3.25$. Then, the 
$y$-coordinate of the vertex is $-0.2(3.25)^2+1.3(3.25)=2.1125$. Hence, because the maximum value of this graph is less than $3$, the moust cannot jump over a $3$ ft. high 
fence.

\section*{60}
(a): For the earth, the function is: $y=-\frac{32}{10000}x^2+x$ while the equation for the moon is: $y=-\frac{5.3}{10000}x^2+x$.\newline
(b): After graphing these two functions, it can be seen that the golf ball travels $312.5$ ft. on earth, and $1886.792$ ft. on the moon.
(c): The ratio of the distances traveled is $\frac{1886.792}{312.5}\approx 6.04$. We can see that $\frac{32}{5.3}\approx 6.04$. Hence, it can be inferred that, if $F(g_1, x_1)=y, 
F(g_2, x_2)=y$ then $\frac{x_1}{x_2}=\frac{g_2}{g_1}$.

\section*{61}
We can set up the equations:\newline $A=wl,\ P=2w+l$ This yields:\newline $l=P-2w$. Hence, we can multiplly both sides by $w$, and then substitue $A$ in:\newline
$A=w(P-2w)=-2w^2+Pw$. Then, we find the vertex of this parabola. The $x$-coordinate of the vertex is $-\frac{P}{2(-2)}=\frac{P}{4}$. Hence, the $y$-coordinate of the vertex 
is $A=-2(\frac{P}{4})^2+P(\frac{P}{4})$\newline$\rightarrow =-2(\frac{P^2}{16})+\frac{P^2}{4}$\newline $\rightarrow \frac{P^2}{-8}+\frac{P^2}{4}$\newline 



\end{document}

